%%%%%%%%%%%%%%%%% DO NOT CHANGE HERE %%%%%%%%%%%%%%%%%%%% 
%%%%%%%%%%%%%%%%%%%%%%%%%%%%%%%%%%%%%%%%%%%%%%%%%%%%%%%%%%{
    \documentclass[twoside,11pt]{article}
    %%%%% PACKAGES %%%%%%
    \usepackage{pgm2016}
    \usepackage{amsmath}
    \usepackage{algorithm}
    \usepackage[noend]{algpseudocode}
    \usepackage{subcaption}
    \usepackage[english]{babel}	
    \usepackage{paralist}	
    \usepackage[lowtilde]{url}
    \usepackage{fixltx2e}
    \usepackage{listings}
    \usepackage{color}
    \usepackage{hyperref}
    \usepackage{natbib}
    \usepackage{auto-pst-pdf}
    \usepackage{pst-all}
    \usepackage{pstricks-add}
    \usepackage{float} % for [H] placement

    %%%%% MACROS %%%%%%
    \algrenewcommand\Return{\State \algorithmicreturn{} }
    \algnewcommand{\LineComment}[1]{\State \(\triangleright\) #1}
    \renewcommand{\thesubfigure}{\roman{subfigure}}
    \definecolor{codegreen}{rgb}{0,0.6,0}
    \definecolor{codegray}{rgb}{0.5,0.5,0.5}
    \definecolor{codepurple}{rgb}{0.58,0,0.82}
    \definecolor{backcolour}{rgb}{0.95,0.95,0.92}
    \lstdefinestyle{mystyle}{
       backgroundcolor=\color{backcolour},  
       commentstyle=\color{codegreen},
       keywordstyle=\color{magenta},
       numberstyle=\tiny\color{codegray},
       stringstyle=\color{codepurple},
       basicstyle=\footnotesize,
       breakatwhitespace=false,        
       breaklines=true,                
       captionpos=b,                    
       keepspaces=true,                
       numbers=left,                    
       numbersep=5pt,                  
       showspaces=false,                
       showstringspaces=false,
       showtabs=false,                  
       tabsize=2
    }
    \lstset{style=mystyle}
%%%%%%%%%%%%%%%%%%%%%%%%%%%%%%%%%%%%%%%%%%%%%%%%%%%%%%%%%% 
%%%%%%%%%%%%%%%%%%%%%%%%%%%%%%%%%%%%%%%%%%%%%%%%%%%%%%%%%% }

%%%%%%%%%%%%%%%%%%%%%%%% CHANGE HERE %%%%%%%%%%%%%%%%%%%% 
%%%%%%%%%%%%%%%%%%%%%%%%%%%%%%%%%%%%%%%%%%%%%%%%%%%%%%%%%% {
\newcommand\course{xx25106}
\newcommand\courseName{CS Seminar}
\newcommand\semester{Fall 2025}
\newcommand\assignmentNumber{01}                             % <-- ASSIGNMENT #
\newcommand\studentName{Abdelrahman I. A. Abushbeka}                  % <-- YOUR NAME
\newcommand\studentEmail{s2520711@u.tsukuba.ac.jp}          % <-- YOUR NAME
\newcommand\studentNumber{202520711}                % <-- STUDENT ID #
%%%%%%%%%%%%%%%%%%%%%%%%%%%%%%%%%%%%%%%%%%%%%%%%%%%%%%%%%% }
%%%%%%%%%%%%%%%%%%%%%%%%%%%%%%%%%%%%%%%%%%%%%%%%%%%%%%%%%%

%%%%%%%%%%%%%%%%% DO NOT CHANGE HERE %%%%%%%%%%%%%%%%%%%% 
%%%%%%%%%%%%%%%%%%%%%%%%%%%%%%%%%%%%%%%%%%%%%%%%%%%%%%%%%%
%{

    \ShortHeadings{University of Tsukuba -  \course ~~ \courseName}{\studentName - \studentNumber}
    \firstpageno{1}
    
    \begin{document}
    
    \title{Difference Subspaces in Change Segmentation: Sentinel-2 Satellite Damage Mapping}
    
    \author{\name \studentName \email \studentEmail \\
    \studentNumber
    \addr
    }
    
    \maketitle
%%%%%%%%%%%%%%%%%%%%%%%%%%%%%%%%%%%%%%%%%%%%%%%%%%%%%%%%%%

%%%%%%%%%%%%%%%%%%%%%%%%%%%%%%%%%%%%%%%%%%%%%%%%%%%%%%%%%% }
\begin{abstract}
\noindent
Accurate and timely assessment of disaster damage from satellite imagery is
crucial for emergency response and long--term reconstruction. This work
describes a two--phase pipeline that combines classical subspace--based change
detection with supervised segmentation on multi--temporal Sentinel--2 data. In
Phase~1 we implement Difference Subspace (DS) change detection on the Onera
Satellite Change Detection (OSCD) benchmark and on the unlabeled MultiSenGE
Sentinel--2 dataset, comparing DS against classical unsupervised baselines:
pixel differencing, change vector analysis (CVA), PCA--based difference
(PCA--diff), Celik's local PCA+K--means method, and IR--MAD. Phase~1 produces
continuous change maps and quantitative OSCD results (AUROC, F1, IoU) that
serve as interpretable priors. In Phase~2 we train U--Net and ResNet--U--Net
models for pixel--wise change segmentation on OSCD, using Sentinel--2 pre/post
bands alone and in combination with DS and PCA--diff change maps as additional
channels and fusion branches. On OSCD, raw Sentinel--2 segmentation achieves
the best IoU/F1, while DS/PCA--diff priors slightly improve AUROC and serve as
valuable explanatory overlays. The code and design are structured so that the
same architecture and priors can be transferred to true damage--labeled
datasets such as xBD, forming the basis of a future Phase~3 for building damage
segmentation.
\end{abstract}

\section{Introduction}

Remote sensing offers a global, repeatable view of urban and rural areas before
and after disasters. However, robust damage assessment from satellite imagery
remains challenging: annotations are expensive, radiometric conditions vary,
and small changes in vegetation or illumination can confound simple difference
metrics. At the same time, classical linear subspace methods provide a
principled way to model spectral manifolds and their changes over time
\cite{FukuiMaki2015,FukuiSubspace}.

The long--term goal of this project is label--efficient \emph{disaster damage
segmentation} from multi--temporal satellite data. Rather than directly
training a deep model on a damage dataset, we decompose the problem into two
complementary phases:

\begin{itemize}
  \item \textbf{Phase 1:} Unsupervised difference--subspace (DS) change
  detection on Sentinel--2, evaluated on the OSCD benchmark and visualized on
  the MultiSenGE Sentinel--2 dataset. Phase~1 produces continuous change maps
  and classical baselines that serve as priors.

  \item \textbf{Phase 2:} Supervised change segmentation on OSCD, using
  Sentinel--2 pre/post bands as inputs and incorporating DS/PCA--diff change
  maps as additional channels or fusion branches. Phase~2 establishes strong
  change--segmentation baselines and tests how much DS/PCA priors help.
\end{itemize}

Both phases are designed with \emph{forward compatibility}: by changing only
the dataset adapter and the number of output classes, the same pipeline can be
ported to damage--labeled datasets such as xBD~\cite{Gupta2019xBD} in a future
Phase~3.

\section{Related Work}

\subsection{Subspace methods and difference subspace}

Subspace methods such as mutual subspace methods and related Grassmannian
learning techniques are widely used for image set recognition and time--series
modeling~\cite{FukuiSubspace}. Fukui and Maki introduced the
\emph{Difference Subspace} (DS) and its generalization for subspace--based
methods, providing a unified framework for separating common and
difference components between two subspaces~\cite{FukuiMaki2015}. DS has been
extended to time--series anomaly detection and higher--order
variants~\cite{FukuiMaki2015,SecondOrderDS}, showing that difference
subspaces can highlight subtle yet systematic changes between signal
subspaces.

In our work, DS is used in a purely spectral sense: each Sentinel--2 image is
represented as a cloud of 13--dimensional spectral vectors, and DS is
constructed between pre-- and post--event spectral subspaces. The resulting
DS projection energy acts as an unsupervised change score; we also compare a
cross--residual variant inspired by illumination--robust formulations.

\subsection{Classical change detection in remote sensing}

Unsupervised change detection in satellite images is a mature field. Simple
methods such as pixel differencing and change vector analysis (CVA) compute
L2 norms between pre/post spectra. PCA--based methods perform dimensionality
reduction on the difference image and use the magnitude of projections onto
leading principal components as change scores, often yielding strong
performance with low computational cost. Celik proposed a local PCA + K--means
change detector that partitions the difference image into local
patches, learns a PCA basis per patch, and clusters projected feature vectors
into change/no--change classes using K--means~\cite{Celik2009}. IR--MAD
iteratively reweights canonical variates between pre and post images to
produce robust multivariate difference maps.

We implement all of these baselines in Phase~1 for comparison with DS.

\subsection{Deep learning for change and damage detection}

Convolutional neural networks (CNNs) have been successfully applied to
multi--temporal image change detection~\cite{Daudt2018}. U--Net
architectures~\cite{Ronneberger2015} are widely adopted for semantic
segmentation, while deep residual networks (ResNets)~\cite{He2016ResNet}
provide strong backbones. For damage assessment, the xBD dataset introduced a
large--scale benchmark of pre/post imagery with multi--level building damage
labels~\cite{Gupta2019xBD}, enabling supervised training of damage
segmentation networks.

Our Phase~2 segmentation experiments are intentionally conservative: we use
standard U--Net and ResNet--U--Net architectures on OSCD binary change masks
to build a clear baseline, and we investigate whether classical DS/PCA priors
add value beyond raw Sentinel--2 bands.

\section{Datasets and Preprocessing}

\subsection{OSCD: Onera Satellite Change Detection}

The Onera Satellite Change Detection (OSCD) dataset consists of 24 Sentinel--2
tiles with pre-- and post--event images (13 spectral bands) and binary change
masks at \SI{10}{\metre}/\SI{20}{\metre}/\SI{60}{\metre}
resolutions~\cite{Daudt2018}. Following standard practice, we treat each city
as a tile with tensors
\(
X_{\text{pre}}, X_{\text{post}} \in \mathbb{R}^{13 \times H \times W}
\)
and a binary change mask
\(Y \in \{0,1\}^{1 \times H \times W}\).

We adopt the official train/test split provided with the dataset (14 train
cities, 10 test cities) and choose a subset of train cities as validation. For
all tiles, we build a valid--pixel mask based on a configurable NODATA value
and a minimum number of valid bands, and we apply global per--band z--score
normalization using statistics computed over the OSCD training set.

\subsection{MultiSenGE Sentinel--2 dataset}

The MultiSenGE dataset provides \(\sim\!8000\) unlabeled
multi--temporal Sentinel--2 patches (256\(\times\)256) over the Grand--Est
region. In Phase~1 we use only the Sentinel--2 L2A bands as an unlabeled
testbed: patches are grouped by tile and spatial index, and date--sorted
pairs (e.g.\ earliest vs.\ latest) are selected for DS visualization. Band
statistics from OSCD are reindexed to the MultiSenGE band order so that
normalization is consistent across datasets.

\section{Phase 1: Difference--Subspace Change Detection}

\subsection{Difference subspace construction}

For a given OSCD tile, we assemble normalized spectral matrices
\(
X_1, X_2 \in \mathbb{R}^{d \times n}
\)
from pre/post images, where \(d=13\) bands and \(n\) is the number of valid
pixels. We fit rank--\(r\) PCA bases using a scikit--learn style PCA:

\begin{align}
\Phi &= \text{PCA}(X_1) \in \mathbb{R}^{d \times r},\\
\Psi &= \text{PCA}(X_2) \in \mathbb{R}^{d \times r},
\end{align}
with either fixed rank or a variance threshold (e.g.\ \(95\%\)).

We implement two DS variants:
\begin{itemize}
  \item \textbf{Residual--stacked DS} (default). Let \(P_\Phi = \Phi \Phi^\top\)
  and \(P_\Psi = \Psi \Psi^\top\). The residual projectors
  \(R_\Phi = I - P_\Phi\), \(R_\Psi = I - P_\Psi\) isolate components outside
  each subspace. The difference subspace basis \(D\) is formed by
  orthonormalizing the concatenation of residual projections:
  \begin{equation}
      D = \operatorname{orth}\left(
      \left[ R_\Psi \Phi,\, R_\Phi \Psi \right]
      \right).
  \end{equation}

  \item \textbf{Eigen--based DS}. Following the generalized DS formulation
  in~\cite{FukuiMaki2015}, we construct
  \(G = P_\Phi + P_\Psi\) and compute its eigenpairs. Eigenvectors associated
  with eigenvalues in \((\varepsilon, 1-\varepsilon)\) span the difference
  subspace, while values near 2 correspond to the common subspace.
\end{itemize}

Empirically, on OSCD both variants produce almost identical scores and
evaluation metrics.

\subsection{DS scores and classical baselines}

Given a pixel pair \((x_1,x_2)\) and basis \(D\), we define two DS scores:
\begin{align}
p_{\text{DS}} &= \lVert D^\top (x_2 - x_1) \rVert_2^2,\\
r_{\text{cross}} &= \lVert R_\Psi x_2 \rVert_2^2 + \lVert R_\Phi x_1 \rVert_2^2.
\end{align}
Phase~1 primarily uses \(p_{\text{DS}}\) as the DS projection score.

We compare DS against:
\begin{itemize}
  \item \textbf{Pixel diff / CVA:} \(s_{\text{pix}} = \lVert x_2 - x_1
  \rVert_2\).
  \item \textbf{PCA--diff:} PCA on the difference image \(X_2 - X_1\), using
  the magnitudes of projections onto leading PCs as scores.
  \item \textbf{Celik PCA+K--means:} local PCA on difference--image patches and
  K--means clustering into change/no--change
  classes~\cite{Celik2009}.
  \item \textbf{IR--MAD:} iteratively reweighted multivariate alteration
  detection, implemented as a stretch baseline.
\end{itemize}

Scores are normalized per tile and converted to binary masks using either
unsupervised Otsu thresholds (per tile) or a train--calibrated global
threshold. We evaluate AUROC on continuous scores and IoU/F1 on binary masks
for train/val/test splits.

\subsection{Phase 1 results}

On the OSCD test split, PCA--diff achieves the highest AUROC (around 0.81),
with F1/IoU slightly above DS and pixel diff. DS projection is competitive
(AUROC \(\sim 0.75\)) and often produces cleaner, less noisy change maps than
direct pixel differencing, while the cross--residual variant and IR--MAD are
weaker. Celik's method performs reasonably but is computationally heavier due
to local PCA and clustering.

Qualitative OSCD figures show that DS and PCA--diff highlight many of the
same regions as the ground truth but also respond to subtle spectral changes
(e.g.\ vegetation, radiometric shifts) not labeled as ``change'' in OSCD. On
MultiSenGE, DS maps provide a broad picture of spectral change across diverse
scenes without any labels.

The key outcome of Phase~1 is thus a set of interpretable, unsupervised
change maps (DS, PCA--diff, pixel diff, etc.) on OSCD and MultiSenGE, together
with a robust evaluation pipeline.

\section{Phase 2: OSCD Change Segmentation with DS/PCA Priors}

\subsection{Segmentation models}

Phase~2 introduces supervised change segmentation on OSCD, using the Phase~1
change maps as priors.

The core baseline is a 2D U--Net~\cite{Ronneberger2015} with 4 encoder/decoder
levels (channels 64--128--256--512) and a final \(1\times 1\) convolution to
produce logits in \(\mathbb{R}^{1 \times H \times W}\). We also implement a
stronger U--Net with a ResNet--34 encoder~\cite{He2016ResNet} and a custom
decoder.

Input feature configurations include:
\begin{itemize}
  \item \textbf{Raw only}: 26 channels, concatenating normalized pre and post
  Sentinel--2 bands.
  \item \textbf{Raw + DS}: 26 raw channels plus DS projection score (27).
  \item \textbf{Raw + PCA}: 26 raw channels plus PCA--diff score (27).
  \item \textbf{Raw + DS + PCA}: 26 raw channels plus both priors (28).
  \item \textbf{Priors only}: DS and PCA--diff maps without raw bands (for
  ablation).
\end{itemize}

In addition to naive concatenation, we design a small
\emph{PriorsFusionUNet}, where raw bands and priors are first processed by
separate \(1\times 1\) convolutions and then fused channel--wise before
entering the U--Net. This allows priors to have their own feature branches.

\subsection{Training and evaluation protocol}

We extract overlapping 256\(\times\)256 patches from OSCD tiles with a
configurable overlap (e.g.\ 64 pixels), and we apply random flips and 90°
rotations to both images and masks. Training uses a combination of
BCEWithLogits and soft Dice loss on valid pixels, with AdamW optimizer and a
cosine learning rate schedule over 50 epochs.

We report mean IoU, mean F1, and AUROC over the OSCD test split, averaging
over all test tiles. For qualitative analysis, Phase~2 includes visualization
scripts that overlay pre/post RGB, ground truth, DS/PCA--diff maps, prediction
probabilities, and binary segmentation masks in combined multi--panel figures.

\subsection{Phase 2 results}

On OSCD, raw Sentinel--2 segmentation is a strong baseline:
\begin{itemize}
  \item A plain U--Net with raw pre/post bands achieves mean IoU
  \(\approx 0.23\), mean F1 \(\approx 0.33\), and AUROC \(\approx 0.87\) on the
  test split.
  \item A ResNet--U--Net with raw bands yields similar IoU/F1 and slightly
  different AUROC, depending on hyperparameters.
\end{itemize}

When DS/PCA--diff change maps are added as extra channels:
\begin{itemize}
  \item Concatenating DS or PCA--diff individually does not improve IoU/F1 and
  often slightly degrades them; combined DS+PCA priors show similar trends.
  \item PriorsFusionUNet improves over naive concatenation in terms of F1 and
  AUROC but still does not surpass the best raw--only models in IoU/F1.
  \item In several configurations, AUROC increases modestly when priors are
  included, suggesting that priors sharpen ranking but also introduce more
  false positives at a fixed threshold.
\end{itemize}

Per--tile visualizations confirm this picture: DS and PCA--diff maps often
highlight broader regions of spectral change than the binary OSCD masks,
including radiometric artefacts and subtle vegetation changes. Networks that
consume priors tend to be more sensitive to these regions, which can be
beneficial for detecting weak changes but harmful for strict IoU/F1 against
the OSCD labels.

\section{Discussion and Future Directions}

The two--phase pipeline provides a coherent story:

\begin{enumerate}
  \item Phase~1 shows that DS and PCA--diff are strong, interpretable
  unsupervised change detectors on Sentinel--2. PCA--diff has the best AUROC
  on OSCD among classical methods, while DS is competitive and often cleaner
  than raw pixel diff.
  \item Phase~2 demonstrates that, at least on OSCD, a well--trained
  segmentation network on raw Sentinel--2 data is hard to beat in IoU/F1 when
  evaluated strictly against the provided change masks.
\end{enumerate}

However, DS/PCA--diff priors remain valuable in two key ways:

\begin{itemize}
  \item \textbf{Explainability.} Combined figures with DS/PCA--diff maps,
  segmentation probabilities, and ground truth help diagnose where the network
  ``follows'' or ``ignores'' unsupervised spectral change signals. This is
  important when moving to real disaster scenarios where labels may be noisy
  or incomplete.
  \item \textbf{Transferability.} The DS machinery is dataset--agnostic. It can
  be recomputed on future damage datasets (e.g.\ xBD~\cite{Gupta2019xBD}) and
  used as priors, pseudo--labels, or loss weighting factors in a Phase~3
  damage--segmentation pipeline.
\end{itemize}

Several natural extensions are left for future work:

\begin{itemize}
  \item Longer training and multiple random seeds to assess variance.
  \item ImageNet pretraining for the ResNet encoder and partial freezing of
  early layers.
  \item Self--supervised or pseudo--label pretraining on MultiSenGE using DS
  or PCA--diff maps as targets.
  \item Alternative uses of priors, such as loss weighting on DS--strong
  regions, or injecting prior features at deeper encoder layers.
  \item Extension from binary change masks (OSCD) to multi--class damage
  levels on xBD or xBD--S12, via a dataset adapter and multi--class output
  heads.
\end{itemize}

\section{Conclusion}

We have described a two--phase Sentinel--2 pipeline that moves from classical
difference subspaces and PCA--based change detection to supervised change
segmentation on OSCD, with DS and PCA--diff maps reused as priors and
explanatory tools. Phase~1 establishes that subspace--based change scores
provide competitive unsupervised performance and interpretable change maps
on both OSCD and MultiSenGE. Phase~2 builds robust U--Net and ResNet--U--Net
segmentation baselines, showing that raw Sentinel--2 inputs are sufficient to
achieve strong IoU/F1 on OSCD, while DS/PCA--diff priors mainly contribute to
AUROC and interpretability.

The codebase and experimental design are structured so that, in a future
Phase~3, the same architecture and DS machinery can be applied to
damage--labeled datasets such as xBD, ultimately targeting fast, interpretable,
and label--efficient damage assessment in real disaster scenarios.

\begin{thebibliography}{99}

\bibitem{FukuiMaki2015}
K.~Fukui and A.~Maki, ``Difference subspace and its generalization for
subspace-based methods,'' \emph{IEEE Transactions on Pattern Analysis and
Machine Intelligence}, vol.~37, no.~11, pp.~2164--2177, 2015.

\bibitem{FukuiSubspace}
K.~Fukui, ``Subspace methods,'' Tech.\ Rep., University of Tsukuba.
(Overview/tutorial document on mutual subspace methods and related
approaches.)

\bibitem{SecondOrderDS}
K.~Fukui and Y.~Kobayashi, ``Second-order difference subspace,'' in
\emph{Proc.\ of a computer vision / pattern recognition workshop},
year omitted here. (Higher-order extension of difference subspace.)

\bibitem{Celik2009}
T.~Celik, ``Unsupervised change detection in satellite images using principal
component analysis and $k$-means clustering,'' \emph{IEEE Geoscience and
Remote Sensing Letters}, vol.~6, no.~4, pp.~772--776, 2009.

\bibitem{Daudt2018}
A.~Daudt, B.~Le~Saux, and A.~Boulch, ``Fully convolutional siamese networks
for change detection,'' in \emph{Proc.\ ICIP}, 2018. (Introduces CNN-based
change detection and the OSCD Sentinel-2 dataset.)

\bibitem{Ronneberger2015}
O.~Ronneberger, P.~Fischer, and T.~Brox, ``U-Net: Convolutional networks for
biomedical image segmentation,'' in \emph{Proc.\ MICCAI}, 2015.

\bibitem{He2016ResNet}
K.~He, X.~Zhang, S.~Ren, and J.~Sun, ``Deep residual learning for image
recognition,'' in \emph{Proc.\ CVPR}, 2016.

\bibitem{Gupta2019xBD}
R.~Gupta \emph{et al.}, ``xBD: A dataset for assessing building damage from
satellite imagery,'' in \emph{Proc.\ CVPR Workshops}, 2019.

\end{thebibliography}

\end{document}